\documentclass{article}

% AI4DD @ NeurIPS 2026 uses the NeurIPS 2026 style. Per the CfP, submit
% anonymized with \usepackage{neurips_2026} (no options). For a camera-ready
% add [final]. Place neurips_2026.sty / .bst next to this file.
\usepackage{neurips_2026}

\usepackage[utf8]{inputenc}
\usepackage[T1]{fontenc}
\usepackage{hyperref}
\usepackage{url}
\usepackage{booktabs}
\usepackage{amsfonts}
\usepackage{nicefrac}
\usepackage{microtype}
\usepackage{graphicx}
\usepackage{amsmath}
\usepackage{xcolor}

% Draft helper: \TODO{...} for gaps to fill (esp. real-data numbers). Remove before submission.
\newcommand{\TODO}[1]{{\color{red}\textbf{[TODO: #1]}}}
\newcommand{\Ed}{\mathcal{E}} % energy distance symbol

\title{Exploration Reduces Mode Averaging in Perturbation-Response
Generative Models: A Distribution-Fidelity View}

% Anonymous for submission.
\author{%
  Anonymous Author(s) \\
  Affiliation \\
  \texttt{email} \\
}

\begin{document}
\maketitle

\begin{abstract}
Generative models of single-cell perturbation response (e.g.\ latent diffusion
models such as scLDM) predict the \emph{distribution} of cell states induced by
a drug or genetic perturbation. Because a perturbation drives a heterogeneous
population---responder and non-responder subpopulations, distinct
transcriptional programs---the target is multimodal, yet models trained with a
per-sample reconstruction loss regress to the mean and smear probability mass
into biologically implausible states between response modes. We ask whether
\emph{Explorative Modeling} (XM), a best-of-$K$ exploration wrapper for
generative models, mitigates this \emph{mode averaging} in the perturbation
setting. Wrapping a conditional flow-matching generator over a single-cell VAE
latent, we evaluate on biologically meaningful axes with an independent
bio-first framework---distribution fidelity (energy distance), differential
expression recovery, and pathway/program recovery---rather than aggregate
$\text{MSE-}\Delta$ alone. Across synthetic Perturb-seq and \TODO{real
Norman/Replogle} data, exploration significantly improves distribution-level
fidelity (energy distance; all paired bootstrap CIs exclude zero), with a
monotonic dose-response in $K$ and the largest gains in the few-step,
fast-inference regime where mode averaging is worst. The gain is a targeted one:
it lands on population/heterogeneity fidelity while mean-level metrics are
approximately unchanged, exposing a precision-versus-calibration trade that
standard benchmarks hide. We frame this as a concrete instance of the
benchmark-to-reality translation gap for generative models in drug discovery.
\end{abstract}

\section{Introduction}
\label{sec:intro}
% Motivation: perturbation-response prediction as generative modeling; why it matters for drug discovery.
Predicting how cells respond to perturbations is central to target
identification, mechanism-of-action, and prioritizing experiments in drug
discovery. Modern approaches cast this as conditional generative modeling in a
learned single-cell latent space \citep{scldm,scvi}.

% The core failure mode.
A single perturbation induces a \emph{heterogeneous} population: responder and
non-responder subpopulations and distinct transcriptional programs make the
conditional target multimodal. Models trained to reconstruct individual cells
under an arbitrary noise-to-data pairing regress toward the conditional mean and
place mass in the empty valley \emph{between} response modes---predicting cells
that do not occur. We call this \emph{mode averaging}.

% The translation-gap angle (AI4DD framing).
This failure is invisible to the aggregate mean-shift metrics
($\text{PCC-}\Delta$, $\text{MSE-}\Delta$) that dominate perturbation
benchmarks: a model can score well on the mean while badly misrepresenting the
population. This is a concrete instance of the benchmark-to-reality translation
gap---strong static-benchmark numbers that need not reflect biological
reliability.

% Contribution.
We study whether \emph{Explorative Modeling} (XM)~\citep{xm}, a modality-agnostic
best-of-$K$ exploration wrapper, mitigates mode averaging for perturbation
response. Our contributions:
\begin{itemize}
  \item We apply XM's exploration to an scLDM-shaped conditional generator
  (flow matching over a single-cell VAE latent), requiring no change to the
  exploration engine (\S\ref{sec:method}).
  \item We evaluate on \emph{biological} axes with an independent framework---
  distribution fidelity, DEG recovery, and pathway/program recovery---with
  bootstrap confidence intervals, rather than aggregate error alone
  (\S\ref{sec:setup}).
  \item We show exploration significantly and monotonically improves
  distribution-level fidelity, most in the fast-inference regime, while
  mean-level metrics are approximately unchanged---a precision-versus-calibration
  trade that standard benchmarks hide (\S\ref{sec:results}).
\end{itemize}

\section{Method}
\label{sec:method}

\paragraph{Perturbation response as conditional generation.}
Let $x\in\mathbb{R}^{G}$ be a cell's expression over $G$ genes and $c$ a
perturbation label. An scLDM-style model learns a conditional generator
$p_\theta(z\mid c)$ in a VAE latent $z=E(x)$, decoding $\hat{x}=D(z)$. We
instantiate the generator as conditional flow matching~\citep{flowmatching,
liprince}: with source noise $z_0\sim\mathcal{N}(0,I)$, data $z_1$, and
$z_t=(1-t)z_0+t z_1$, a velocity network $v_\theta(z_t,t,c)$ is trained to
regress the target $z_1-z_0$.

\paragraph{Mode averaging.}
At a noised point $z_t$ the optimal regression target is the average velocity
over all data consistent with $z_t$; for a multimodal conditional this average
points between modes, and few-step integration lands in the low-density valley.
This is precisely the population-level error that aggregate metrics miss.

\paragraph{Explorative Modeling.}
XM~\citep{xm} adds a third pretraining axis via best-of-$K$ exploration: at each
step it draws $K$ candidate latent variables (here, source noises), scores each
by the model's loss against the fixed data target, and back-propagates only
through the best candidate (Forward XM). Each update thus commits to a single
mode rather than averaging across many. Concretely the training loss for a data
point becomes $\min_{k\in\{1,\dots,K\}} \ell\big(v_\theta, z_1, z_0^{(k)}\big)$;
$K{=}1$ recovers the ordinary generator. We use the reference implementation's
exploration engine unchanged, wrapping the flow loss in its selection routine.

\section{Experimental setup}
\label{sec:setup}

\paragraph{Models.}
A conditional velocity MLP over the latent, trained as baseline ($K{=}1$) and
with exploration ($K\in\{2,4,8\}$), identical except for $K$. We consider three
latent spaces to test robustness: standardized $\log$-expression, a Gaussian VAE
latent, and a scVI-style negative-binomial (NB) count-decoder latent
\citep{scvi}---the count-likelihood generative model used in practice.

\paragraph{Data.}
(i) A synthetic Perturb-seq with raw counts, gene knockdowns, and an injected
responder/non-responder split per perturbation, giving genuine downstream DEGs
\emph{and} the multimodal response that induces mode averaging. (ii) Real
Perturb-seq---\TODO{Norman 2019 \citep{norman}; Replogle 2022 \citep{replogle}}
---loaded through the bio-perturbations dataset adapters. The task is
distribution recovery on \emph{seen} perturbations (cell-level held-out truth);
the model conditions on a learned per-perturbation embedding.

\paragraph{Evaluation.}
We score predictions with an independent bio-first evaluator across four
families: aggregate mean-shift ($\text{PCC-}\Delta$, $\text{MSE-}\Delta$);
\emph{distribution} fidelity via the energy distance $\Ed$~\citep{energydist}
between predicted and held-out perturbed populations; \emph{DEG recovery}
(directional recall@$k$, top-$k$ Jaccard, effect-size correlation), with
truth-side DEGs from sample-matched pseudobulk DESeq2~\citep{deseq2}; and
\emph{pathway/program recovery} (up/down Jaccard and rank correlation) against
known gene programs. We report bootstrap $95\%$ CIs over
$(\text{seed}\times\text{perturbation})$ units and paired CIs on the
XM-vs-baseline difference.

\paragraph{Baselines.}
Identity (control unchanged; the floor), MeanShift (control $+$ mean training
$\Delta$), and scGen-style latent vector arithmetic~\citep{scgen}. \TODO{Add an
established perturbation model, e.g.\ GEARS \citep{gears} or CPA
\citep{cpa}, on real data.}

\section{Results}
\label{sec:results}

\paragraph{Exploration reduces mode averaging.}
Table~\ref{tab:main} reports the main comparison. Exploration lowers the energy
distance to the held-out perturbed population---the metric that captures whether
the predicted \emph{population} reproduces response heterogeneity---while the
trivial and scGen baselines lag. \TODO{Fill from real-data run
\texttt{scldm\_ablation\_<dataset>\_nbvae.json}.}

\begin{table}[t]
\centering
\caption{Main results (\emph{synthetic Perturb-seq; replace with real
Norman/Replogle}), NB-VAE latent, mean over 3 seeds; $\downarrow$/$\uparrow$ =
lower/higher better. $\Ed$ = distribution/heterogeneity fidelity; DEG rec.@20 =
directional recall; DEG $r$ = effect-size Pearson on true DEGs; Path.\ $\rho$ =
program rank correlation (down-programs).}
\label{tab:main}
\small
\begin{tabular}{lcccccc}
\toprule
Model & $\Ed\downarrow$ & PCC-$\Delta\uparrow$ & MSE-$\Delta\downarrow$ &
DEG rec.@20$\uparrow$ & DEG $r\uparrow$ & Path.\ $\rho\uparrow$ \\
\midrule
Identity (floor)      & 8.32 & $-0.00$ & 11.38 & 0.24 & 0.06 & $-0.03$ \\
MeanShift             & 9.17 & 0.32 & 11.59 & 0.32 & 0.75 & --- \\
scGen (latent arith.) & 8.26 & \TODO{} & \TODO{} & \TODO{} & \TODO{} & \TODO{} \\
mini-scLDM, $K{=}1$   & 8.41 & \textbf{0.86} & \textbf{4.99} & \textbf{0.51} & 0.991 & 0.29 \\
mini-scLDM, XM $K{=}4$& \textbf{8.18} & 0.85 & 5.58 & \textbf{0.51} & \textbf{0.992} & \textbf{0.48} \\
\bottomrule
\end{tabular}
\end{table}

\paragraph{Biological signal: DEG and pathway recovery.}
% Whether the improvement is (a) real biology, not just a distance number, and (b) not bought at the cost of DEG recovery.
The distribution-level gain is not bought by degrading downstream biology
(Table~\ref{tab:main}). Relative to the $K{=}1$ generator, exploration leaves DEG
directional recall unchanged ($0.51\!\to\!0.51$), slightly improves top-$20$
DEG Jaccard ($0.42\!\to\!0.44$) and the effect-size correlation on true DEGs
($0.991\!\to\!0.992$), and \emph{improves} program recovery, most for
down-regulated programs (rank correlation $0.29\!\to\!0.48$). Thus on synthetic
data the improved population fidelity comes essentially for free on the DEG axis
and with a modest gain on pathway/program recovery; the only cost is a small
rise in $\text{MSE-}\Delta$ (\S``precision--calibration trade''). \TODO{Confirm
this pattern holds on real data---the central biological claim of the paper.}

\paragraph{Ablation: exploration budget and sampling steps.}
Figure~\ref{fig:ablation} sweeps $K$ and sampling steps. The improvement is
significant at every setting (paired CIs exclude zero), monotonic in $K$, and
largest at few sampling steps---the fast-inference regime where mode averaging
is worst.

\begin{figure}[t]
\centering
% copy the PNG into paper/ before building, or adjust the path
\includegraphics[width=0.72\linewidth]{figs/ablation.png}
\caption{Energy distance vs.\ exploration budget $K$ for $2/4/10$ sampling
steps (VAE latent, 5 seeds, $95\%$ CI bands); dashed lines are baselines.
Lower is better. \emph{Synthetic; regenerate on real data with
\texttt{scldm\_ablation.py --plot}.}}
\label{fig:ablation}
\end{figure}

\paragraph{The precision--calibration trade.}
The gain is targeted: exploration improves $\Ed$ (population fidelity) while
mean-level $\text{PCC-}\Delta$ and DEG recovery are approximately unchanged and
$\text{MSE-}\Delta$ can tick up---consistent with Forward XM sharpening modes at
some cost to distributional calibration. \TODO{State the real-data direction and
magnitude; discuss whether Reverse XM narrows it.}

\paragraph{Latent-space robustness.}
The effect is consistent across $\log$-expression, Gaussian-VAE, and
NB-VAE latents (Appendix~\TODO{}), i.e.\ it survives the encode--decode
round-trip through a true count decoder---the space an scLDM actually models.

\section{Discussion and limitations}
\label{sec:discussion}
% Translation-gap framing + honest limitations.
Our results argue that for generative perturbation models, \emph{which} metric
one optimizes matters: exploration buys population/heterogeneity fidelity that
aggregate error hides, directly relevant to the benchmark-to-reality gap this
workshop targets. Limitations: (i) \TODO{extent of real-data evidence};
(ii) pseudo-replicate pseudobulk underestimates biological variance;
(iii) a single generator class and a compact VAE rather than a full scLDM;
(iv) the precision--calibration trade means XM as-is is not a strict improvement
on every axis. Future work: Reverse XM for calibration, stronger baselines
(GEARS/CPA), and larger real panels.

\section*{Reproducibility}
Code, synthetic generator, XM wrapper, three latent spaces, the real-dataset
loader path, and the ablation/CIs are released (\texttt{experiments/perturbation\_xm/});
evaluation uses the open bio-perturbations framework~\citep{bioperturbations}.

\bibliographystyle{plainnat}
\bibliography{references}

\end{document}
